\chapter{Index Theorems and Riemann--Roch}\label{index-theorems}
\chapterstatus{complete}

\section{Introduction}\label{index-theorems:section-introduction}

An elliptic operator on a compact manifold has a finite-dimensional kernel and cokernel (Theorem~\ref{elliptic-operators:theorem-fredholm}), and the difference of their dimensions, the \emph{index}, is remarkably stable: it does
not change under deformations, and the Atiyah--Singer index theorem computes it from characteristic classes. For the three classical elliptic complexes the index is a familiar invariant: the Euler characteristic, the
signature, and the holomorphic Euler characteristic $\chi(X, E) = \sum_q(-1)^q\dim H^q(X, E)$. The corresponding cases of the index theorem are the Gauss--Bonnet--Chern theorem, Hirzebruch's signature theorem and the
Hirzebruch--Riemann--Roch theorem. This chapter identifies the three indices with complete proofs, states the index theorem and its three cases, proves Hirzebruch--Riemann--Roch directly for line bundles on curves and for
all line bundles on projective spaces, and draws the integrality consequences (Noether's formula, divisibility of Chern numbers) that make these theorems useful in algebraic geometry. References: \cite{hirzebruch-topological},
\cite{atiyah-singer-1968-1}, \cite{atiyah-singer-1968-3}, \cite{lawson-michelsohn}, \cite{griffiths-harris}.

\noindent\textbf{Prerequisites.} Chapter~\ref{hodge-theorem} (Harmonic Forms and the Hodge Theorem), Chapter~\ref{dolbeault} (Dolbeault Cohomology), Chapter~\ref{characteristic-classes} (Characteristic Classes),
Chapter~\ref{curves} (Algebraic Curves), Chapter~\ref{coherent-cohomology} (Cohomology of Coherent Sheaves) and Chapter~\ref{cycle-class} (The Cycle Class Map).

\noindent\textbf{Conventions.} Manifolds are compact, and characteristic classes are taken in de Rham cohomology (or rational cohomology), identified by the de Rham theorem. For a real vector bundle $V$ we write
$V_\CC = V \otimes_\RR \CC$. For a class $a$ in the cohomology of a compact oriented manifold, $\int_Ma$ is the integral of its top-degree component.

\section{The index of an elliptic operator}\label{index-theorems:section-index}

\begin{definition}\label{index-theorems:definition-index}
The \emph{index} of an elliptic differential operator $P \colon \Gamma(M, E) \to \Gamma(M, F)$ on a compact manifold is $\operatorname{ind}P = \dim\ker P - \dim\ker P^*$, where $P^*$ is the formal adjoint
(Definition~\ref{elliptic-operators:definition-formal-adjoint}).
\end{definition}

By Theorem~\ref{elliptic-operators:theorem-fredholm}, $\ker P^*$ is isomorphic to the cokernel of $P$ on smooth sections, so $\operatorname{ind}P = \dim\ker P - \dim\operatorname{coker}P$, and it is the Fredholm index of
$P \colon H^{s+m}(E) \to H^s(F)$ for every $s$. As the Fredholm index is locally constant on the space of Fredholm operators (Chapter~\ref{functional-analysis}), $\operatorname{ind}P$ is unchanged when $P$ is deformed through a
continuous family of elliptic operators of the same order; in particular it depends only on the principal symbol up to homotopy through invertible symbols.

\begin{proposition}[Euler characteristic operator]\label{index-theorems:proposition-euler-operator}
For a compact oriented Riemannian manifold, $d + \delta \colon \Omega^{\mathrm{even}}(M) \to \Omega^{\mathrm{odd}}(M)$ is elliptic with index $\chi(M)$.
\end{proposition}

\begin{proof}
This is Corollary~\ref{hodge-theorem:corollary-euler-characteristic}; ellipticity holds because $(d + \delta)^2 = \Delta$ has principal symbol $|\xi|^2$, so the symbol of $d + \delta$ is injective for $\xi \neq 0$, and it maps between bundles of
the same rank ($2^{n-1}$).
\end{proof}

\begin{proposition}[Dolbeault operator]\label{index-theorems:proposition-dolbeault-operator}
Let $X$ be a compact complex manifold with a Hermitian metric and $E$ a holomorphic vector bundle with a Hermitian metric. The operator
$\bar\partial_E + \bar\partial_E^* \colon \mathcal{A}^{0,\mathrm{even}}(X, E) \to \mathcal{A}^{0,\mathrm{odd}}(X, E)$ is elliptic, and its index is $\chi(X, E) = \sum_q(-1)^q\dim H^q(X, E)$.
\end{proposition}

\begin{proof}
Let $D = \bar\partial_E + \bar\partial_E^*$ on $\mathcal{A}^{0,\bullet}(X, E)$. As $\bar\partial_E^2 = 0$, $D^2 = \bar\partial_E\bar\partial_E^* + \bar\partial_E^*\bar\partial_E = \Box_E$, which is elliptic with principal symbol $\frac12|\xi|^2$
(Theorem~\ref{dolbeault:theorem-dbar-hodge}); so $D$ is elliptic, and it is formally self-adjoint, so the formal adjoint of its even part is its odd part. If $D\alpha = 0$, then $\Box_E\alpha = 0$; conversely
$(\Box_E\alpha, \alpha) = \|D\alpha\|^2$, so $\ker D = \ker\Box_E = \bigoplus_q\mathcal{H}^{0,q}(X, E)$, split into even and odd $q$. By Theorem~\ref{dolbeault:theorem-dbar-hodge}, $\mathcal{H}^{0,q}(X, E) \cong H^q(X, E)$, and the index is
$\sum_q(-1)^q\dim H^q(X, E)$.
\end{proof}

\begin{proposition}[Signature operator]\label{index-theorems:proposition-signature-operator}
Let $M$ be a compact oriented Riemannian manifold of dimension $4m$, and $\tau$ the operator on $\Omega^\bullet(M) \otimes \CC$ equal to $i^{k(k-1)+2m}{*}$ on $k$-forms. Then $\tau^2 = 1$, $\tau$ anticommutes with $d + \delta$,
and $d + \delta \colon \Omega^+ \to \Omega^-$, between the $\pm1$-eigenspaces of $\tau$, is elliptic with index the signature $\sigma(M)$ (Definition~\ref{poincare-duality:definition-intersection-form}).
\end{proposition}

\begin{proof}[Sketch of proof]
We omit the verification of $\tau^2 = 1$ and $\tau(d + \delta) = -(d + \delta)\tau$, which follow from $** = (-1)^{k(n-k)}$ and $\delta = -{*}d{*}$ in even dimension (Proposition~\ref{riemannian-geometry:proposition-adjoint}) by
tracking signs \cite[Chapter 3]{lawson-michelsohn}. Granting these, $d + \delta$ maps $\Omega^+$ to $\Omega^-$, the kernel of $d + \delta$ on $\Omega^\pm$ is the space $\mathcal{H}^\pm$ of harmonic forms in $\Omega^\pm$, and the index is
$\dim\mathcal{H}^+ - \dim\mathcal{H}^-$, exactly as in Proposition~\ref{index-theorems:proposition-euler-operator}. As $\tau$ commutes with $\Delta$ and maps harmonic $k$-forms to harmonic $(4m - k)$-forms, for $k < 2m$ the space
$\mathcal{H}^k \oplus \mathcal{H}^{4m-k}$ is $\tau$-stable, and $\alpha \mapsto \alpha \pm \tau\alpha$ identifies $\mathcal{H}^k$ with its $\pm1$-eigenspaces, which therefore have equal dimensions and cancel in the index. On
$\mathcal{H}^{2m}$, $\tau = i^{2m(2m-1)+2m}{*} = *$, and for a real harmonic $2m$-form $\alpha$ with $*\alpha = \pm\alpha$, $\int_M\alpha \wedge \alpha = \pm\int_M\alpha \wedge *\alpha = \pm\|\alpha\|^2$. So the intersection form is positive definite on
the self-dual and negative definite on the anti-self-dual harmonic forms, which are orthogonal for it, and by the Hodge theorem the index $\dim\mathcal{H}^{2m}_+ - \dim\mathcal{H}^{2m}_-$ is the signature.
\end{proof}

\section{The Atiyah--Singer index theorem}\label{index-theorems:section-atiyah-singer}

\begin{definition}\label{index-theorems:definition-genera}
For a complex vector bundle with Chern roots $x_1, \ldots, x_r$ and a real vector bundle $V$ whose complexification has Chern roots $\pm y_1, \ldots, \pm y_s$ (and $0$ if the rank is odd), define
\[
\mathrm{td}(E) = \prod_i\frac{x_i}{1 - e^{-x_i}}, \qquad L(V) = \prod_j\frac{y_j}{\tanh y_j}, \qquad \hat A(V) = \prod_j\frac{y_j/2}{\sinh(y_j/2)}.
\]
$L$ and $\hat A$ are even in each $y_j$, hence polynomials in the Pontryagin classes: $L = 1 + \frac13p_1 + \frac{1}{45}(7p_2 - p_1^2) + \cdots$ and $\hat A = 1 - \frac{1}{24}p_1 + \cdots$. For a compact oriented $4m$-manifold, the
\emph{$L$-genus} is $\int_ML(TM)$; for a compact complex manifold, the \emph{Todd genus} is $\int_X\mathrm{td}(T_X)$.
\end{definition}

\begin{theorem}[Atiyah--Singer]\label{index-theorems:theorem-atiyah-singer}
The index of an elliptic operator on a compact manifold depends only on the class of its principal symbol in the compactly supported K-theory of the tangent bundle, and it is given by a universal cohomological formula in
the Chern character of that class and the Todd class of $TM_\CC$. For the three classical operators the formula gives:
\begin{enumerate}
\item (Gauss--Bonnet--Chern) $\chi(M) = \int_Me(TM)$ for a compact oriented manifold;
\item (Hirzebruch signature theorem) $\sigma(M) = \int_ML(TM)$ for a compact oriented $4m$-manifold;
\item (Hirzebruch--Riemann--Roch) $\chi(X, E) = \int_X\mathrm{ch}(E)\,\mathrm{td}(T_X)$ for a holomorphic vector bundle $E$ on a compact complex manifold $X$, where $T_X$ is the holomorphic tangent bundle.
\end{enumerate}
\end{theorem}

This is \cite{atiyah-singer-1968-1} and \cite{atiyah-singer-1968-3}; for the formulation and the heat equation proof see \cite{lawson-michelsohn}. The signature theorem and Riemann--Roch for projective manifolds were proved
earlier by Hirzebruch, using Thom's cobordism theory \cite{hirzebruch-topological}. Case (1) is also Theorem~\ref{characteristic-classes:theorem-euler}(3). For the Dirac operator of a spin manifold the index is the
$\hat A$-genus, which is therefore an integer on spin manifolds.

\begin{remark}\label{index-theorems:remark-grr}
For a projective manifold $X$ and an algebraic vector bundle $E$, case (3) also follows from the Grothendieck--Riemann--Roch theorem (Theorem~\ref{intersection-theory:theorem-grr}): it gives
$\chi(X, E) = \deg(\mathrm{ch}(E)\,\mathrm{td}(T_X))$ in the Chow ring, and the cycle class map is a ring homomorphism compatible with Chern classes and degrees (Theorem~\ref{cycle-class:theorem-cycle-class}), so the right-hand
side is $\int_X\mathrm{ch}(E)\,\mathrm{td}(T_X)$; by GAGA the algebraic and analytic Euler characteristics agree (Chapter~\ref{gaga}).
\end{remark}

\section{Riemann--Roch: direct verifications}\label{index-theorems:section-verifications}

\begin{proposition}\label{index-theorems:proposition-rr-curves}
Hirzebruch--Riemann--Roch holds for line bundles on a nonsingular projective curve $X$ of genus $g$: $\chi(X, \mathcal{O}(D)) = \int_X\mathrm{ch}(\mathcal{O}(D))\,\mathrm{td}(T_X) = \deg D + 1 - g$.
\end{proposition}

\begin{proof}
The left side is $l(D) - h^1(\mathcal{O}(D)) = l(D) - l(K - D)$ by Serre duality (Theorem~\ref{coherent-cohomology:theorem-serre-duality}), which is $\deg D + 1 - g$ by the Riemann--Roch theorem
(Theorem~\ref{curves:theorem-riemann-roch}). On the right, $\mathrm{ch}(\mathcal{O}(D)) = 1 + c_1(\mathcal{O}(D))$ and $\mathrm{td}(T_X) = 1 + \frac12c_1(T_X)$ in $H^\bullet(X)$, as $H^{>2}(X) = 0$; so the integrand of degree $2$ is
$c_1(\mathcal{O}(D)) + \frac12c_1(T_X)$. Now $\int_Xc_1(\mathcal{O}(D)) = \deg D$, because $c_1(\mathcal{O}(D))$ is the cycle class of $D$ and the cycle class map is compatible with degrees
(Theorem~\ref{cycle-class:theorem-cycle-class}); and $c_1(T_X) = -c_1(K_X)$ with $\deg K = 2g - 2$ (Corollary~\ref{curves:corollary-riemann-roch-consequences}). So the right side is $\deg D - (g - 1)$.
\end{proof}

\begin{lemma}\label{index-theorems:lemma-todd-projective}
On $\PP^n$ with hyperplane class $h$, $\mathrm{td}(T_{\PP^n}) = \big(h/(1 - e^{-h})\big)^{n+1}$ and $\mathrm{ch}(\mathcal{O}(k)) = e^{kh}$.
\end{lemma}

\begin{proof}
The Euler sequence $0 \to \mathcal{O} \to \mathcal{O}(1)^{\oplus(n+1)} \to T_{\PP^n} \to 0$ (Example~\ref{vector-bundles:example-euler-sequence}) and the multiplicativity of the Todd class in exact sequences
(Proposition~\ref{characteristic-classes:proposition-chern-character}) give $\mathrm{td}(T)\,\mathrm{td}(\mathcal{O}) = \mathrm{td}(\mathcal{O}(1))^{n+1}$, with $\mathrm{td}(\mathcal{O}) = 1$ and $\mathrm{td}(\mathcal{O}(1)) = h/(1 - e^{-h})$,
as $c_1(\mathcal{O}(1)) = h$ (Example~\ref{characteristic-classes:example-tautological}). Also $\mathrm{ch}(\mathcal{O}(k)) = e^{c_1(\mathcal{O}(k))} = e^{kh}$.
\end{proof}

\begin{proposition}\label{index-theorems:proposition-rr-projective}
For all $n \geq 1$ and $k \in \ZZ$,
\[
\chi(\PP^n, \mathcal{O}(k)) = \int_{\PP^n}\mathrm{ch}(\mathcal{O}(k))\,\mathrm{td}(T_{\PP^n}) = \binom{n + k}{n} = \frac{(k + 1)(k + 2)\cdots(k + n)}{n!}.
\]
\end{proposition}

\begin{proof}
\emph{The left side.} By Theorem~\ref{coherent-cohomology:theorem-projective-space}, $H^q(\PP^n, \mathcal{O}(k)) = 0$ for $0 < q < n$; $h^0(\mathcal{O}(k))$ is the number of monomials of degree $k$ in $n + 1$ variables, which is
$\binom{n+k}{n}$ for $k \geq 0$ and $0$ for $k < 0$; and $h^n(\mathcal{O}(k))$ is the number of Laurent monomials $x_0^{a_0}\cdots x_n^{a_n}$ with all $a_i \leq -1$ and $\sum a_i = k$, which is the number of monomials of degree $-k - n - 1$,
that is $\binom{-k-1}{n}$ for $k \leq -n - 1$ and $0$ otherwise. For $k \geq 0$ the polynomial $P(k) = (k + 1)\cdots(k + n)/n!$ equals $\binom{n+k}{n}$; for $-n \leq k \leq -1$ it vanishes, as does $\chi$; and for $k \leq -n - 1$,
$P(k) = (-1)^n(-k - 1)(-k - 2)\cdots(-k - n)/n! = (-1)^n\binom{-k-1}{n} = (-1)^nh^n(\mathcal{O}(k))$. So $\chi(\PP^n, \mathcal{O}(k)) = P(k)$ in all cases.

\emph{The right side.} By Lemma~\ref{index-theorems:lemma-todd-projective} and $\int_{\PP^n}h^n = 1$, the integral is the coefficient of $h^n$ in $e^{kh}(h/(1 - e^{-h}))^{n+1}$, a power series convergent near $0$, that is, the residue at
$0$ of the meromorphic function $e^{kh}(1 - e^{-h})^{-n-1}$:
\[
\frac{1}{2\pi i}\oint_{|h| = \varepsilon}\frac{e^{kh}\,dh}{(1 - e^{-h})^{n+1}}.
\]
The substitution $u = 1 - e^{-h}$ is biholomorphic near $0$ with $u(0) = 0$, so it maps the small circle to a loop winding once around $0$; since $e^{-h} = 1 - u$, $e^{kh} = (1 - u)^{-k}$ and $dh = du/(1 - u)$, the integral becomes
$\frac{1}{2\pi i}\oint(1 - u)^{-k-1}u^{-n-1}\,du$, the coefficient of $u^n$ in $(1 - u)^{-k-1} = \sum_j\binom{k + j}{j}u^j$ (Chapter~\ref{complex-analysis}). This coefficient is $\binom{k+n}{n} = P(k)$, read as a polynomial in $k$.
\end{proof}

\begin{corollary}[Noether's formula]\label{index-theorems:corollary-noether}
For a compact complex surface $S$, $\chi(S, \mathcal{O}_S) = \frac{1}{12}\int_S(c_1^2 + c_2)$, where $c_i = c_i(T_S)$. Consequently $\int_S(c_1^2 + c_2)$ is divisible by $12$, and for a line bundle $L$ with $l = c_1(L)$,
$\chi(S, L) = \chi(S, \mathcal{O}_S) + \frac12\int_Sl(l + c_1)$.
\end{corollary}

\begin{proof}
By Definition~\ref{index-theorems:definition-genera}, the degree-$4$ part of the Todd class is $\frac{1}{12}(c_1^2 + c_2)$ (Definition~\ref{characteristic-classes:definition-chern-character}), and Theorem~\ref{index-theorems:theorem-atiyah-singer}(3) with
$E = \mathcal{O}_S$ gives the formula; the left side is an integer. For $L$, the degree-$4$ part of $(1 + l + \frac12l^2)(1 + \frac12c_1 + \frac{1}{12}(c_1^2 + c_2))$ is $\frac12l^2 + \frac12lc_1 + \frac{1}{12}(c_1^2 + c_2)$.
\end{proof}

\begin{lemma}\label{index-theorems:lemma-pontryagin-complex}
For a complex vector bundle $E$, $(E_\RR)_\CC \cong E \oplus \bar E$, and hence $p_1(E_\RR) = c_1(E)^2 - 2c_2(E)$.
\end{lemma}

\begin{proof}
The complex structure $J$ of $E$ extends $\CC$-linearly to $(E_\RR)_\CC$, where it has eigenvalues $\pm i$; the $+i$-eigenbundle $\{v - iJv\}$ is isomorphic to $E$ by $v \mapsto \frac12(v - iJv)$, and the $-i$-eigenbundle to $\bar E$. By
Definition~\ref{characteristic-classes:definition-pontryagin} and the Whitney formula, $p_1(E_\RR) = -c_2(E \oplus \bar E) = -(c_2(E) + c_1(E)c_1(\bar E) + c_2(\bar E)) = c_1^2 - 2c_2$, using $c_1(\bar E) = -c_1(E)$ and $c_2(\bar E) = c_2(E)$
(Proposition~\ref{characteristic-classes:proposition-chern-properties}(4)).
\end{proof}

\begin{example}\label{index-theorems:example-surfaces}
\begin{enumerate}
\item $\PP^2$: $c(T) = (1 + h)^3$, so $c_1^2 = 9$, $c_2 = 3$ and Noether's formula gives $\chi(\mathcal{O}) = 12/12 = 1$, as it must. By Lemma~\ref{index-theorems:lemma-pontryagin-complex}, $p_1 = 9 - 6 = 3$, and the signature theorem
$\sigma = \frac13\int p_1 = 1$ agrees with the intersection form $(1)$ on $H^2(\PP^2) = \ZZ h$.
\item $\PP^1 \times \PP^1$: $c_1 = 2a + 2b$ and $c_2 = 4ab$, where $a$, $b$ are the pullbacks of the point classes, with $a^2 = b^2 = 0$ and $\int ab = 1$. So $c_1^2 = 8$, $c_2 = 4$, $\chi(\mathcal{O}) = 1$, and $p_1 = 8 - 8 = 0$, in agreement with the
hyperbolic intersection form on $H^2$, of signature $0$.
\item A K3 surface has $c_1 = 0$ and $\chi(\mathcal{O}) = h^0 - h^1 + h^2 = 1 - 0 + 1 = 2$ (Chapter~\ref{calabi-yau}), so Noether's formula gives $\int c_2 = 24$, the topological Euler characteristic; then $p_1 = -48$ and
$\sigma = -16$. As $b_2 = 22$, the intersection form has $3$ positive and $19$ negative eigenvalues.
\end{enumerate}
\end{example}

\begin{counterexample}\label{index-theorems:counterexample-integrality}
The characteristic numbers appearing in the index theorems are a priori rational, and their integrality is a real constraint. There is no compact complex surface with $\int c_1^2 = \int c_2 = 1$, since $12 \nmid 2$ would contradict
Noether's formula. By the signature theorem, $\int_Mp_1(TM) = 3\sigma(M)$ is divisible by $3$ for every compact oriented $4$-manifold $M$; so no such manifold has first Pontryagin number $1$, although nothing in the formal properties
of characteristic classes excludes it.
\end{counterexample}

\section{Exercises}\label{index-theorems:section-exercises}

\begin{exercise}\label{index-theorems:exercise-quadric}
Using Hirzebruch--Riemann--Roch on $\PP^1 \times \PP^1$, show that $\chi(\mathcal{O}(p, q)) = (p + 1)(q + 1)$ for the line bundle of bidegree $(p, q)$, and check it for $p, q \geq 0$ by counting bihomogeneous polynomials.
\end{exercise}

\begin{solution}
With $l = pa + qb$ and $c_1 = 2a + 2b$, Corollary~\ref{index-theorems:corollary-noether} gives $\chi = 1 + \frac12\int(pa + qb)((p + 2)a + (q + 2)b) = 1 + \frac12(p(q + 2) + q(p + 2)) = pq + p + q + 1$. For $p, q \geq 0$ the higher cohomology
vanishes (by the K\"unneth formula and Theorem~\ref{coherent-cohomology:theorem-projective-space}) and the sections are the polynomials of bidegree $(p, q)$, of which there are $(p + 1)(q + 1)$.
\end{solution}

\begin{exercise}\label{index-theorems:exercise-todd-genus}
Show that the Todd genus of $\PP^n$ is $1$, and that the Todd genus of a compact complex manifold equals $\chi(X, \mathcal{O}_X) = \sum_q(-1)^qh^{0,q}(X)$.
\end{exercise}

\begin{solution}
The second statement is Theorem~\ref{index-theorems:theorem-atiyah-singer}(3) with $E = \mathcal{O}_X$, since $\mathrm{ch}(\mathcal{O}_X) = 1$. For $\PP^n$, Proposition~\ref{index-theorems:proposition-rr-projective} with $k = 0$ gives
$\binom{n}{n} = 1$.
\end{solution}

\begin{exercise}\label{index-theorems:exercise-hypersurface-euler}
Let $S \subseteq \PP^3$ be a smooth surface of degree $d$. Using the adjunction formula $c_1(T_S) = (4 - d)h$ and Noether's formula, show that $\int_Sc_2(T_S) = d^3 - 4d^2 + 6d$, and check that this gives $24$ for a quartic.
\end{exercise}

\begin{solution}
From the exact sequence $0 \to \mathcal{O}_{\PP^3}(-d) \to \mathcal{O}_{\PP^3} \to \mathcal{O}_S \to 0$, additivity of $\chi$ and Proposition~\ref{index-theorems:proposition-rr-projective},
$\chi(\mathcal{O}_S) = 1 - \frac{(1 - d)(2 - d)(3 - d)}{6} = 1 + \frac{(d - 1)(d - 2)(d - 3)}{6}$. As
$\int_Sh^2 = d$, $\int c_1^2 = (4 - d)^2d$. Noether's formula gives $\int c_2 = 12\chi(\mathcal{O}_S) - (4 - d)^2d = 12 + 2(d - 1)(d - 2)(d - 3) - d(4 - d)^2 = d^3 - 4d^2 + 6d$. For $d = 4$: $64 - 64 + 24 = 24$.
\end{solution}
